% ============================================================
% 第 8 章 函数与导数
% ============================================================
\chapter{8}{函数与导数}

\subchapter{1}{基础知识}

\begin{definition}[函数单调性]
设函数 $f(x)$ 的定义域为 $D$，区间 $I \subseteq D$。若对任意 $x_1, x_2 \in I$，当 $x_1 < x_2$ 时都有
$f(x_1) < f(x_2)$，则称 $f(x)$ 在区间 $I$ 上单调递增。
\end{definition}

\begin{theorem}[求导基本公式]
$f(x) = x^n$ 的导数为 $f'(x) = nx^{n-1}$；$f(x)=\ln x$ 的导数为 $f'(x)=\dfrac{1}{x}$；
$f(x)=e^x$ 的导数为 $f'(x)=e^x$。
\end{theorem}

\subchapter{2}{例题精讲}

\begin{example}[利用导数求极值]
设 $f(x)=x^3-3x+1$，求 $f(x)$ 的极值。
\end{example}

求导得 $f'(x)=3x^2-3=3(x+1)(x-1)$，令 $f'(x)=0$ 得 $x=\pm 1$。
当 $x<-1$ 时 $f'(x)>0$；当 $-1<x<1$ 时 $f'(x)<0$；当 $x>1$ 时 $f'(x)>0$。
故 $f(x)$ 在 $x=-1$ 处取得极大值 $f(-1)=3$，在 $x=1$ 处取得极小值 $f(1)=-1$。

\begin{example}[参数取值范围]
已知 $f(x)=\ln x-ax$ 在 $[1,+\infty)$ 上单调递减，求 $a$ 的取值范围。
\end{example}

由题意 $f'(x)=\dfrac{1}{x}-a\le 0$ 在 $[1,+\infty)$ 上恒成立，即 $a\ge \dfrac{1}{x}$ 对任意 $x\ge 1$ 恒成立。
由于 $\dfrac{1}{x}$ 在 $[1,+\infty)$ 上的最大值为 $1$，故 $a\ge 1$。

\subchapter{3}{巩固练习}

\begin{exercise}[求极值]
设 $f(x)=x^3-3x+1$，求 $f(x)$ 的极值。
\end{exercise}

\begin{solution}
$f'(x)=3(x+1)(x-1)$，极大值 $f(-1)=3$，极小值 $f(1)=-1$。
\end{solution}

\begin{exercise}[恒成立问题]
已知 $f(x)=\ln x-ax$ 在 $(0,+\infty)$ 上单调递减，求 $a$ 的取值范围。
\end{exercise}

\begin{exercise}[指数不等式]
证明：$e^x\ge x+1$ 对任意实数 $x$ 成立。
\end{exercise}

\begin{solution}
设 $g(x)=e^x-x-1$，则 $g'(x)=e^x-1$。$g(x)$ 在 $(-\infty,0)$ 上递减，在 $(0,+\infty)$ 上递增，
故 $g(x)\ge g(0)=0$，即 $e^x\ge x+1$。
\end{solution}
